\section*{第 \ref{ch:b1:findim} 章 --- 有限维}

\begin{solution}{exo:b1:findim:1}
\begin{enumerate}
  \item $z = -x - y$: $F = \{(x, y, -x-y)\} =
        \operatorname{Vect}\bigl((1,0,-1), (0,1,-1)\bigr)$；这两个
        向量线性无关（由坐标即可看出）：$\dim F = 2$。
  \item $G = \{(x, x, 2t, t)\} =
        \operatorname{Vect}\bigl((1,1,0,0), (0,0,2,1)\bigr)$：线性无关，
        $\dim G = 2$.
  \item $P(1) = P'(1) = 0$ 意味着 $(X-1)^2 \mid P$
        （由\cref{prop:b1:poly:multiplicity}），$P = (X-1)^2(aX + b)$。
        基为 $\bigl((X-1)^2, X(X-1)^2\bigr)$，维数为 $2$。
\end{enumerate}
\end{solution}

\begin{solution}{exo:b1:findim:2}
$(3,5,7) = (1,1,1) + 2(1,2,3)$ 与 $(0,1,2) = (1,2,3) - (1,1,1)$：
两者都是前两个向量的线性组合，而前两个向量线性无关（彼此不成比例）。
秩为 $2$；张成的子空间的基为：
$\bigl((1,1,1), (1,2,3)\bigr)$。
\end{solution}

\begin{solution}{exo:b1:findim:3}
不妨添加 $e_1 = (1,0,0,0)$ 与 $e_3 = (0,0,1,0)$。向量族
$\bigl((1,1,0,0), (0,0,1,1), e_1, e_3\bigr)$ 线性无关：考察
零线性组合 $\alpha(1,1,0,0) + \beta(0,0,1,1) + \gamma e_1 + \delta
e_3 = 0$ 写作 $(\alpha + \gamma, \alpha, \beta + \delta, \beta) =
0$，于是 $\alpha = \beta = 0$，进而 $\gamma = \delta = 0$。四个
线性无关向量，维数为 $4$，故为基
        （见\cref{prop:b1:findim:twoofthree}）。
\end{solution}

\begin{solution}{exo:b1:findim:4}
次数 $0, 1, 2, 3$ 两两不同：线性无关
        （由\cref{prop:b1:vspaces:freecriteria}）可知，四个向量
维数为 $4$，故构成基。对 $X^3$ 作向下展开：
\[
X(X-1)(X-2) = X^3 - 3X^2 + 2X,
\qquad
X(X-1) = X^2 - X,
\]
于是 $X^3 - X(X-1)(X-2) = 3X^2 - 2X$；又 $3X^2 - 2X = 3(X^2 - X) + X
= 3\,X(X-1) + X$。由此得
\[
X^3 = X(X-1)(X-2) + 3\,X(X-1) + 1\cdot X + 0\cdot 1 :
\]
坐标为 $(0, 1, 3, 1)$，相对于基 $\bigl(1, X, X(X-1),
X(X-1)(X-2)\bigr)$。（这些数实际上是斯特林数。）
\end{solution}

\begin{solution}{exo:b1:findim:5}
由格拉斯曼公式 $\dim(F \cap G) = 4 + 5 - \dim(F + G)$，且 $F + G$ 是
$E$ 中包含 $G$ 的子空间：$5 \leq \dim(F+G) \leq 7$。由此
$\dim(F \cap G) \in \{2, 3, 4\}$。在 $\R^7$ 中构造实例，
以标准基 $(e_1, \dots, e_7)$，取 $G =
\operatorname{Vect}(e_1, \dots, e_5)$:
\begin{itemize}
  \item $F = \operatorname{Vect}(e_1, e_2, e_6, e_7)$: $F + G = \R^7$,
        交空间为 $\operatorname{Vect}(e_1, e_2)$，维数为 $2$；
  \item $F = \operatorname{Vect}(e_1, e_2, e_3, e_6)$：交空间
        的维数为 $3$；
  \item $F = \operatorname{Vect}(e_1, e_2, e_3, e_4) \subseteq G$:
        维数为 $4$。
\end{itemize}
\end{solution}

\begin{solution}{exo:b1:findim:6}
$H$ 是子空间（与\cref{exo:b1:vspaces:1}~(4) 完全相同）。由
因式定理（见\cref{thm:b1:poly:factor}），$P \in H \iff P =
(X-1)Q$，且 $\deg Q \leq n - 1$。映射 $Q \mapsto (X-1)Q$ 为
从 $\R_{n-1}[X]$ 到 $H$ 上的线性双射，于是 $H$ 的一组基为
\[
\bigl((X-1),\ (X-1)X,\ (X-1)X^2,\ \dots,\ (X-1)X^{n-1}\bigr),
\qquad \dim H = n .
\]
一个补子空间为 $\operatorname{Vect}(1)$（常值多项式）。事实上
$H \cap \operatorname{Vect}(1) = \{0\}$（非零常数在
$1$ 处不为零），且维数之和为 $n + 1$：由格拉斯曼公式，$H
\oplus \operatorname{Vect}(1) = \R_n[X]$.
\end{solution}

\begin{solution}{exo:b1:findim:7}
删除：删去 $u_p$ 后，张成子空间只能缩小；并且它
至多缩小一维：将 $u_p$ 加回较小张成子空间的基，
即得较大张成子空间的一个生成组，其中至多多出
一个向量。对称地，添加向量 $v$：新的
张成子空间包含原来的，至多多出一个生成元，因此其
维数为 $r$（若 $v$ 已在张成子空间中）或 $r + 1$
（否则，由\cref{prop:b1:vspaces:freecriteria}~(2)，
线性无关组可扩充为基）。两者合起来，即得“秩是
$1$-利普希茨”的结论。
\end{solution}

\begin{solution}{exo:b1:findim:8}
沿严格递增的链，维数严格增加
（$F_i \subseteq F_{i+1}$，$F_i \neq F_{i+1}$，且
由\cref{thm:b1:findim:subspaces}，维数相等将迫使
空间相等）。于是 $\dim F_1 < \dim F_2 < \dots < \dim F_k$ 是
$\intint{0}{n}$ 中的严格递增整数序列：至多有
$n + 1$ 个值，$k \leq n + 1$。$\R^n$ 中的极大链为：
\[
\{0\} \subsetneq \operatorname{Vect}(e_1) \subsetneq
\operatorname{Vect}(e_1, e_2) \subsetneq \dots \subsetneq \R^n ,
\]
长度恰为 $n + 1$。
\end{solution}

\begin{solution}{exo:b1:findim:9}
$F + G$ 严格包含 $F$（因为 $G \not\subseteq F$），于是
$\dim(F + G) = n$，由格拉斯曼公式得 $\dim(F \cap G) = (n-1) +
(n-1) - n = n - 2$.

一般断言：对 $k$ 作归纳，$I_k$ 是 $k$
个超平面的交，满足 $\dim I_k \geq n - k$。$k = 1$ 时显然成立。归纳步：$I_{k+1}
= I_k \cap H_{k+1}$，并在 $E$ 中应用格拉斯曼公式：
\[
\dim(I_k \cap H_{k+1}) = \dim I_k + (n - 1) - \dim(I_k + H_{k+1})
\geq \dim I_k + (n-1) - n \geq n - k - 1 . \qedhere
\]
\end{solution}

\begin{solution}{exo:b1:findim:10}
\begin{enumerate}
  \item 条件 $u_{n+2} - 3u_{n+1} + 2u_n = 0$ 是线性的，
        且零数列满足它：$E$ 是子空间。由
        归纳法，$u_0$ 与 $u_1$ 决定每个 $u_n$，且这种依赖
        是线性的（每一步都是前两个值的
        线性组合）；反之，每一对 $(a, b)$
        恰好对应 $E$ 中唯一的数列（由
        递推式定义 $u_n$）。同\cref{ex:b1:findim:dims}，$E$ 以
        双射且线性的方式由 $(u_0, u_1) \in
        \R^2$ 参数化：$\dim E = 2$。
  \item 常数列：$3\cdot1 - 2\cdot1 = 1$。等比数列：$3\cdot
        2^{n+1} - 2\cdot 2^n = (6 - 2)2^n = 2^{n+2}$。两者都属于
        $E$。线性无关：由 $a\cdot 1 + b\cdot 2^n = 0$ 对所有 $n$
        成立，取 $n = 0$ 与 $n = 1$ 得 $a + b = 0$，$a + 2b = 0$，
        故 $a = b = 0$。维数为 $2$ 的空间中，两个线性无关的向量
        构成一组基（\cref{prop:b1:findim:twoofthree}）。
  \item 解 $a + b = 0$，$a + 2b = 1$：得 $b = 1$，$a = -1$，故
        $u_n = 2^n - 1$（即梅森数列）。
\end{enumerate}
\end{solution}

\begin{solution}{exo:b1:findim:11}
\begin{enumerate}
  \item 由格拉斯曼公式，$\dim(F \cap G) = \dim F + \dim G - \dim(F +
        G) \geq \dim F + \dim G - n > 0$，故 $F \cap G \neq
        \{0\}$。当 $n = 100$ 时：$51 + 50 - 100 = 1 > 0$，
        故交集含一条直线。
  \item 若 $F \cap H = \{0\}$，则和为直和，且 $\dim F +
        (n - 1) = \dim(F \oplus H) \leq n$，故 $\dim F \leq 1$。
        （反之，不含于 $H$ 的直线确实满足
        此条件：超平面几乎不“挡”任何东西。）
\end{enumerate}
\end{solution}

\begin{solution}{exo:b1:findim:12}
对 $d = \dim F = \dim G$ 作向下归纳，从 $d = n$ 到 $d =
0$。若 $d = n$，则 $F = G = E$，$S = \{0\}$ 满足要求。假设
该命题对维数为 $d + 1 \leq n$ 的子空间对成立，
并设 $\dim F = \dim G = d < n$。

若 $F = G$：取 $S$ 为 $F$ 的任一补子空间
（\cref{thm:b1:findim:subspaces}）。若 $F \neq G$：两者都是真子空间，
由\cref{exo:b1:vspaces:12}，存在 $x \notin F \cup G$。
和 $F' = F \oplus \operatorname{Vect}(x)$ 与 $G' = G \oplus
\operatorname{Vect}(x)$ 均为直和（$x \notin F$，$x \notin G$），
且维数为 $d + 1$。由归纳假设，它们
有公共补子空间 $S'$：$E = F' \oplus S' = G' \oplus
S'$。令 $S = \operatorname{Vect}(x) \oplus S'$ --- 直和，
因为 $S' \cap \operatorname{Vect}(x) \subseteq S' \cap F' =
\{0\}$。

于是 $F + S = F + \operatorname{Vect}(x) + S' = F' + S' = E$，且
$F \cap S = \{0\}$：若 $f = \lambda x + s'$，其中 $f \in F$，$s'
\in S'$，则 $s' = f - \lambda x \in F' \cap S' = \{0\}$，故 $f
= \lambda x$，于是 $\lambda = 0$（因 $x \notin F$），从而 $f = 0$。
因此 $E = F \oplus S$，同理 $E = G \oplus S$。
\end{solution}

\begin{solution}{pb:b1:findim:1}
\textbf{1.} $\Q$ 含 $0 \neq 1$，对加法、
乘法及取相反数封闭，且每个非零有理数都有有理逆元，
因此是域（\cref{def:b1:structures:field}）。
张成、线性无关与基的定义，以及
第 18--19 章的证明只用到域中的向量加法、分配律与标量
运算，从不涉及绝对值、序或
极限。于是 $\R$ 连同自身的加法，以及
由 $\Q \times \R \to \R$ 给出的数乘，构成 $\Q$ 上的向量空间，所有
定理均适用。除以标量只在
\cref{thm:b1:findim:exchange}的证明中出现一次：要“解出 $g_k$”，需
除以其非零系数。

\textbf{2.} 若 $a + b\sqrt2 = 0$，其中 $a, b \in \Q$ 不全
为零：$b \neq 0$ 将给出 $\sqrt2 = -a/b \in \Q$，与
$\sqrt2$ 的无理性相矛盾（见\cref{exo:b1:arith:7}，取 $p =
2$）；于是 $b = 0$，再得 $a = 0$，即在 $\Q$ 上线性无关。在 $\R$ 上
关系 $\sqrt2\cdot 1 + (-1)\cdot\sqrt2 = 0$ 非平凡，
故线性相关。因此 $\dim_\Q \Q(\sqrt2) = 2$，且坐标
$(a, b)$ 唯一。

\textbf{3.} $(a + b\sqrt2)(c + d\sqrt2) = (ac + 2bd) + (ad +
bc)\sqrt2 \in \Q(\sqrt2)$。于是
\[
\sigma(uv) = (ac + 2bd) - (ad + bc)\sqrt2
= (a - b\sqrt2)(c - d\sqrt2) = \sigma(u)\sigma(v),
\]
故 $N(uv) = uv\,\sigma(uv) = u\sigma(u)\,v\sigma(v) = N(u)N(v)$。
若 $N(u) = a^2 - 2b^2 = 0$ 且 $u \neq 0$：$b \neq 0$ 将给出
$(a/b)^2 = 2$，即 $2$ 有有理平方根；故 $b = 0$，进而 $a
= 0$，从而 $u = 0$，矛盾。因此 $N(u) \neq 0$，当 $u \neq
0$ 时。

\textbf{4.} $u \cdot \dfrac{\sigma(u)}{N(u)} = \dfrac{N(u)}{N(u)}
= 1$，且 $\sigma(u)/N(u) \in \Q(\sqrt2)$：每个非零元素
在 $\Q(\sqrt2)$ 中都可逆，因此它是
$\R$ 的子域。例如 $N(3 + 2\sqrt2) = 9 - 8 = 1$，于是
\[
\frac1{3 + 2\sqrt2} = 3 - 2\sqrt2 ;
\qquad
N(1 + \sqrt2) = -1, \quad
\frac1{1 + \sqrt2} = \sqrt2 - 1 .
\]

\textbf{5.} 若 $\sqrt6 = p/q$，则 $6q^2 = p^2$；
$2$ 的指数在左边为 $1 + 2v_2(q)$，为奇数；在右边为 $2v_2(p)$，为偶数，
矛盾。今设 $\sqrt3 = a + b\sqrt2$，其中
$a, b \in \Q$。两边平方得 $3 = a^2 + 2b^2 + 2ab\sqrt2$。若 $ab
\neq 0$，则 $\sqrt2 = (3 - a^2 - 2b^2)/(2ab) \in \Q$，
矛盾。若 $b = 0$，则 $\sqrt3 = a \in \Q$，与
\cref{exo:b1:arith:7}（$p = 3$）相矛盾。若 $a = 0$，则 $\sqrt3 = b\sqrt2$，
两边乘以 $\sqrt2$，得 $\sqrt6 = 2b \in \Q$，矛盾。
因此 $\sqrt3 \notin \Q(\sqrt2)$。

\textbf{6.} 基向量之间的乘积为
\[
\sqrt2\,\sqrt3 = \sqrt6,\quad
\sqrt2\,\sqrt6 = 2\sqrt3,\quad
\sqrt3\,\sqrt6 = 3\sqrt2,\quad
(\sqrt2)^2 = 2,\ (\sqrt3)^2 = 3,\ (\sqrt6)^2 = 6,
\]
它们均属于 $M$。$M$ 中两元素之积经双线性性展开为
这些元素的有理线性组合，因此 $M$ 对乘法
封闭。

\textbf{7.} 设 $x + y\sqrt3 = 0$，其中 $x, y \in K = \Q(\sqrt2)$。
若 $y \neq 0$，则 $\sqrt3 = -x/y \in K$（第 4 题：$K$ 是
域），与第 5 题矛盾。于是 $y = 0$，进而 $x = 0$：$(1,
\sqrt3)$ 在 $K$ 上线性无关。它生成：$K + K\sqrt3 =
\{(a + b\sqrt2) + (c + d\sqrt2)\sqrt3\} = \operatorname{Vect}_\Q
(1, \sqrt2, \sqrt3, \sqrt6) = M$。于是 $\dim_K M = 2$。

\textbf{8.} 设关系式 $a + b\sqrt2 + c\sqrt3 + d\sqrt6 = 0$
可改写为 $(a + b\sqrt2) + (c + d\sqrt2)\sqrt3 = 0$，其系数
属于 $K$；由第 7 题知两者皆为零，再由第 2 题
得 $a = b = 0$ 且 $c = d = 0$。因此该族线性无关，且 $\dim_\Q M
= 4$。

\textbf{9.} 对每个 $x \in M$，写 $x = \sum_j y_j f_j$，其中 $y_j
\in K$（$M$ 在 $K$ 上的基），并且每个 $y_j = \sum_i
\lambda_{ij} e_i$，其中 $\lambda_{ij} \in \Q$（$K$ 在 $\Q$ 上
的基），代入得 $x = \sum_{i,j} \lambda_{ij}\, e_i f_j$：
这些乘积生成 $M$（视为 $\Q$ 上的向量空间）。

\textbf{10.} 设 $\sum_{i,j} \lambda_{ij}\, e_i f_j = 0$，其中
$\lambda_{ij} \in \Q$。重新合并得 $\sum_j \bigl(\sum_i \lambda_{ij}
e_i\bigr) f_j = 0$，且括号内的和属于 $K$。由
$(f_j)$ 在 $K$ 上线性无关，得 $\sum_i \lambda_{ij} e_i = 0$
对每个 $j$ 成立；再由 $(e_i)$ 在 $\Q$ 上线性无关，
得 $\lambda_{ij} = 0$，对所有 $i, j$ 成立。

\textbf{11.} 由第 9、10 题，$(e_i f_j)_{i \leq m,\, j
\leq n}$ 是 $M$ 在 $\Q$ 上的基，共 $mn$ 个元素：
\[
\dim_\Q M = m\,n = \dim_\Q K \cdot \dim_K M .
\]
验算：$\dim_\Q K = 2$ 与 $\dim_K M = 2$（由第 2、7 题）
给出 $\dim_\Q M = 4$，与第 8 题相符。

\textbf{12.} 向量 $u a_1, \dots, u a_d$ 属于 $A$
（由稳定性）。它们线性无关：若 $\sum_i \lambda_i\, u a_i = 0$，
则 $u \sum_i \lambda_i a_i = 0$（在 $\R$ 中），而 $u \neq 0$ 迫使
$\sum_i \lambda_i a_i = 0$，从而 $\lambda_i = 0$（$a_i$ 构成
基）。于是 $(u a_1, \dots, u a_d)$ 是 $d$ 个
属于 $A$ 的线性无关向量，且 $\dim_\Q A = d$，故为基
（\cref{prop:b1:findim:twoofthree}）。特别地，$1 \in A$
可分解为 $1 = \sum_i \mu_i\, u a_i = u\,v$，其中 $v = \sum_i
\mu_i a_i \in A$，即 $u$ 的逆元属于 $A$。这重新
得到第 4 题的结论（$A = \Q(\sqrt2)$），并且一举证明了 $M$
是 $\R$ 的子域（结合第 6 题）。

\textbf{13.} 这 $d + 1$ 个向量 $1, x, x^2, \dots, x^{d}$ 都
属于 $A$（由乘积的封闭性）；$A$ 中线性无关族至多含有
$d$ 个向量（\cref{prop:b1:findim:twoofthree}），因此它们
线性相关：存在不全为零的有理数 $\lambda_0, \dots, \lambda_d$，
使得 $\sum_k \lambda_k x^k = 0$。多项式 $P
= \sum_k \lambda_k X^k$ 非零，系数为有理数，
次数 $\leq d$，且 $P(x) = 0$。

\textbf{14.} $s^2 = 2 + 2\sqrt6 + 3 = 5 + 2\sqrt6$，坐标为
$(5, 0, 0, 2)$。于是
\[
s^3 = s\,s^2 = (\sqrt2 + \sqrt3)(5 + 2\sqrt6)
= 5\sqrt2 + 2\sqrt{12} + 5\sqrt3 + 2\sqrt{18}
= 11\sqrt2 + 9\sqrt3 ,
\]
坐标为 $(0, 11, 9, 0)$（利用 $\sqrt{12} = 2\sqrt3$，
$\sqrt{18} = 3\sqrt2$）。最后，$s^4 = (s^2)^2 = 25 + 20\sqrt6 +
24 = 49 + 20\sqrt6$，坐标为 $(49, 0, 0, 20)$。

\textbf{15.} 设 $a + bs + cs^2 + ds^3 = 0$。读出四个
坐标（基为 $(1, \sqrt2, \sqrt3, \sqrt6)$）：
\[
a + 5c = 0,\qquad b + 11d = 0,\qquad b + 9d = 0,\qquad 2c = 0 .
\]
于是 $c = 0$，进而 $a = 0$；将中间两式相减，得 $2d =
0$，进而 $b = 0$：$(1, s, s^2, s^3)$ 线性无关。四个线性无关向量
位于 $M$ 中（维数为 $4$）并构成基，故
$\operatorname{Vect}_\Q(1, s, s^2, s^3) = M$。由第 14 题，
$s^3 - 9s = 2\sqrt2$ 与 $11s - s^3 = 2\sqrt3$：
\[
\sqrt2 = \frac{s^3 - 9s}{2},
\qquad
\sqrt3 = \frac{11s - s^3}{2} .
\]

\textbf{16.} $s^4 - 10s^2 + 1 = (49 + 20\sqrt6) - 10(5 +
2\sqrt6) + 1 = 0$。首一、次数 $\leq 3$ 的
多项式若在 $s$ 处取零，则给出 $(1, s, s^2, s^3)$ 的非平凡
零组合，与第 15 题矛盾：$X^4 - 10X^2 + 1$
的次数最低。其根为 $X^2 = 5 \pm 2\sqrt6 = (\sqrt3 \pm
\sqrt2)^2$，故四个实根为 $\pm(\sqrt3 + \sqrt2)$ 与
$\pm(\sqrt3 - \sqrt2)$，即 $\pm\sqrt2 \pm \sqrt3$。

\textbf{17.} 由 $s^4 - 10s^2 + 1 = 0$ 得 $s\,(10s - s^3) = 1$，于是
\[
\frac1s = 10s - s^3 = 10(\sqrt2 + \sqrt3) - (11\sqrt2 + 9\sqrt3)
= \sqrt3 - \sqrt2 ,
\]
与 $(\sqrt3 + \sqrt2)(\sqrt3 - \sqrt2) = 1$ 相符。若 $a +
b\sqrt2 + c\sqrt3 + d\sqrt6 = r \in \Q$，则 $(a - r) + b\sqrt2
+ c\sqrt3 + d\sqrt6 = 0$，由第 8 题的线性无关性知 $b = c
= d = 0$（且 $a = r$）。对 $\sqrt2 + \sqrt3 + \sqrt6$，
其坐标 $(0, 1, 1, 1)$ 不具上述形式，因此该数是无理数 ---
一个比 \cref{exo:b1:reals:7} 更强的命题，
且无需任何平方技巧即可得到。

\textbf{18.} 设 $p/q$（既约）为 $X^3 - 2$ 的有理根，
满足 $p^3 = 2q^3$，由
\cref{exo:b1:poly:5} 的判别准则得 $p \mid 2$，$q \mid 1$：候选
$\pm1, \pm2$，其立方为 $\pm1, \pm8 \neq 2$。无有理
根；特别地 $t = 2^{1/3} \notin \Q$，故 $(1, t)$ 在
$\Q$ 上线性无关（如第 2 题）。

\textbf{19.} 设 $(1, t, t^2)$ 线性相关：存在非零 $P \in
\Q[X]$，$\deg P \leq 2$ 且 $P(t) = 0$；取这样的 $P$，使其次数
\emph{最低}（记为 $d \geq 1$）。由第 18 题，$d \neq 1$，故 $d
= 2$。由带余除法（\cref{thm:b1:poly:division}）：$X^3 - 2
= PQ + R$，其中 $Q, R \in \Q[X]$，$\deg R < 2$。代入 $t$，
得 $0 = P(t)Q(t) + R(t) = R(t)$，故 $R$ 在 $t$ 处为零；由
$d$ 的最小性知 $R = 0$。于是 $X^3 - 2 = PQ$，其中 $\deg Q = 1$：
$Q$ 的有理根即为 $X^3 - 2$ 的有理根，
与第 18 题矛盾。因此 $(1, t, t^2)$ 线性无关，且
$\dim_\Q A = 3$。稳定性：$t^3 = 2$ 将 $1, t, t^2$ 的任意乘积
化为它们的线性组合（$t\cdot t^2 = 2$，$t^2\cdot
t^2 = 2t$）；$A$ 含 $1$，由第 12 题，$A$ 是
$\R$ 的子域。

\textbf{20.} 设 $t \in M$。则 $t^2 \in M$（由稳定性），故
$A \subseteq M$，且 $M$ 是域 $A$ 上的向量空间：
其公理即 $\R$ 中运算法则的限制。它
在 $A$ 上是有限维的（$\Q$ 上的有限生成族
自然也在 $A \supseteq \Q$ 上生成）。塔律
用于 $\Q \subseteq A \subseteq M$，得
\[
4 = \dim_\Q M = \dim_\Q A \cdot \dim_A M = 3\,\dim_A M ,
\]
矛盾：$3$ 不整除 $4$。故 $2^{1/3} \notin M$：
即 $1, \sqrt2, \sqrt3, \sqrt6$ 的有理线性组合都不等于
$2^{1/3}$。

\textbf{21.} $\Q(\sqrt2) \subseteq M$，故 $t \notin \Q(\sqrt2)$。
若 $t + \sqrt2 = r \in \Q$，则 $t = r - \sqrt2 \in \Q(\sqrt2)$，
矛盾 --- $2^{1/3} + \sqrt2$ 为无理数。若 $t = a +
b\sqrt3$，则 $t \in \operatorname{Vect}_\Q(1, \sqrt3) \subseteq
M$，又一矛盾。

\textbf{22.} $B$ 是域 $A$ 上的向量空间
（数乘限制），由于 $\dim_\Q B$ 有限，
故其亦为有限维。对 $\Q \subseteq A \subseteq B$ 用塔律，得
$\dim_\Q B = \dim_\Q A \cdot \dim_A B$：左侧因子整除。
因此 $M$ 的子域在 $\Q$ 上的维数为 $1$、$2$ 或 $4$：
维数为 $1$ 的即 $\Q$ 本身，维数为 $4$ 的即 $M$。

\textbf{23.} 设 $F \subseteq M$ 为满足 $\dim_\Q F =
2$，取 $w \in F \setminus \Q$，则 $(1, w)$ 线性无关，
从而为 $F$ 的基。由第 13 题（$d = 2$），$w^2 = \alpha + \beta w$，其中
$\alpha, \beta$ 为有理数。取 $v = w - \beta/2 \in F
\setminus \Q$，计算得
\[
v^2 = w^2 - \beta w + \frac{\beta^2}4 = \alpha + \frac{\beta^2}4
\;\in\; \Q,
\]
且 $F = \operatorname{Vect}(1, v)$。今设 $v = a + b\sqrt2 +
c\sqrt3 + d\sqrt6$，并展开：
\[
v^2 = (a^2 + 2b^2 + 3c^2 + 6d^2)
+ (2ab + 6cd)\sqrt2 + (2ac + 4bd)\sqrt3 + (2ad + 2bc)\sqrt6 .
\]
三个无理分量为零：$ab + 3cd = ac + 2bd =
ad + bc = 0$。若 $a \neq 0$，则 $b = -3cd/a$，第二个
条件化为 $c(a^2 - 6d^2) = 0$；若 $a^2 = 6d^2$ 且 $d
\neq 0$ 将使 $\sqrt6 = \abs{a/d}$ 为有理数，故须 $c = 0$
或 $d = 0$，两种情形下其余条件均迫使 $b =
c = d = 0$，即 $v \in \Q$，与假设矛盾。故 $a = 0$，且
条件为 $3cd = 2bd = bc = 0$，即 $b, c, d$ 中至多一个
非零。故 $v$ 为 $\sqrt2$、$\sqrt3$
或 $\sqrt6$ 的有理倍数，于是
\[
F \in \bigl\{\Q(\sqrt2),\ \Q(\sqrt3),\ \Q(\sqrt6)\bigr\} ,
\]
这三者确为 $2$ 维子域（稳定性如第 6 题，
逆元如第 12 题）：恰有三个二次
子域。

\textbf{24.} $(1 + \sqrt2 + \sqrt3)(1 + \sqrt2 - \sqrt3) = (1 +
\sqrt2)^2 - 3 = 2\sqrt2$，故
\[
\frac1{1 + \sqrt2 + \sqrt3}
= \frac{1 + \sqrt2 - \sqrt3}{2\sqrt2}
= \frac{\sqrt2 + 2 - \sqrt6}{4}
= \frac12 + \frac14\sqrt2 + 0\cdot\sqrt3 - \frac14\sqrt6 .
\]
（验算：$(1 + \sqrt2 + \sqrt3)(2 + \sqrt2 - \sqrt6) = 4$，
展开后即得。）

\textbf{25.} (i) $\Q$ 上的维数为 $\R$ 的每个子域赋予一个\emph{整数}，
而塔律使这些整数
沿包含关系相乘：维数的行为类似算术
不变量，而整除性约束化为不可能性
证明。(ii) 有限维性迫使每个元素满足一个
次数至多为维数的非零有理系数多项式方程，
即有限性意味着代数性。(iii) “三取二”规则把
“乘以 $u$ 将基映为
极大规模的线性无关组”转化为满射性，
由此得到 $u^{-1}$，无需任何公式：逆元来自计数。(iv) 塔律
禁止 $3$ 维域含于一个
$4$ 维域之中，这正是 $2^{1/3}$ 不能
由 $\sqrt2$ 与 $\sqrt3$ 得到的原因。第 III 部分的定理是
\emph{戴德金塔律}，伽罗瓦理论的开端，
将在本科三年级分册中展开。
\end{solution}
